% STATUS: written 2026-11-05 (Dixmier 1968 pp. 209--211 and 241--242 read;
% Belov-Kanel--Kontsevich 2007 pp. 1--2 read).
\section{The Weyl algebra and the problems of Dixmier}\label{sec:polynomial-weyl:weyl}

Section~\ref{sec:polynomial-weyl:polynomial} closed with the Jacobian
conjecture. This section follows it into the noncommutative world: the
polynomial ring becomes the algebra of Definition~\ref{def:pw:weyl} below. % claim-ok: narration of the section plan, the definition carries the citation
The conjecture becomes the first problem of Dixmier \cite[p.~241]{Dixmier1968}.
Dixmier named the algebra in 1968, after a suggestion of Segal
\cite[p.~209]{Dixmier1968}.

\begin{definition}[{\cite[p.~209]{Dixmier1968}}]\label{def:pw:weyl}
Let $k$ be a commutative field of characteristic $0$ and $n\geq1$ an
integer. The \emph{$n$-th Weyl algebra} $A_{n}(k)$, or $A_{n}$, is the
unital associative $k$-algebra with $2n$ generators
$p_{1},q_{1},\dots,p_{n},q_{n}$ and the relations
\[
[p_{i},q_{i}]=1,\qquad [p_{i},p_{j}]=[q_{i},q_{j}]=[p_{i},q_{j}]=0
\quad(i\neq j),
\]
and $A_{0}(k)=k$.
\end{definition}

The ordered products $p_{1}^{i_{1}}\cdots p_{n}^{i_{n}}q_{1}^{j_{1}}\cdots
q_{n}^{j_{n}}$ form a basis of $A_{n}$ over $k$, and $A_{n}$ is a simple
domain with center $k$, left and right Noetherian
\cite[p.~210]{Dixmier1968}. The algebra is a model of the derivative:
$A_{n}$ is the algebra of polynomial differential operators on affine
$n$-space, with $q_{i}$ acting as multiplication by $x_{i}$ and $p_{i}$ as
$\partial/\partial x_{i}$ \cite[p.~1]{BelovKanelKontsevich2007}. The
derivations of the algebra are all inner \cite[p.~210]{Dixmier1968}.

\begin{proposition}[{\cite[p.~210]{Dixmier1968}}]\label{prop:pw:weyl-inner}
Every derivation of the Weyl algebra $A_{n}$ is inner.
\end{proposition}

The automorphisms are harder: the exponential of the locally nilpotent
derivation $\operatorname{ad}(\lambda q^{n+1})$ gives a well-defined
automorphism of $A_{1}$ \cite[p.~210]{Dixmier1968}. Dixmier posed six
problems at the end of the paper \cite[p.~241]{Dixmier1968}, and the
first is still open.

\begin{remark}[{\cite[Probl\`eme 11.1, p.~241]{Dixmier1968}}]
\label{rem:pw:dixmier-problem}
Dixmier asked whether every endomorphism of $A_{1}$ is an automorphism:
if $P$ and $Q$ in $A_{1}$ satisfy $[P,Q]=1$, do they generate $A_{1}$?
The question in rank $n$ is now called the Dixmier conjecture
\cite[p.~1]{BelovKanelKontsevich2007}.
\end{remark}

Belov-Kanel and Kontsevich tied the conjecture back to the polynomial
thread of Section~\ref{sec:polynomial-weyl:polynomial}. Their theorem
makes the two conjectures one question \cite[p.~2]{BelovKanelKontsevich2007}.

\begin{theorem}[{\cite[Theorem~1, p.~2]{BelovKanelKontsevich2007}}]
\label{thm:pw:bkk}
The Jacobian conjecture in dimension $2n$ implies the Dixmier conjecture
in rank $n$. In particular the stable conjectures $JC_{\infty}$ and
$DC_{\infty}$ are equivalent.
\end{theorem}

The implication $DC_{n}\Rightarrow JC_{n}$ was already known
\cite[p.~1]{BelovKanelKontsevich2007}. Theorem~\ref{thm:pw:bkk} is the
reverse, so the two threads of the chapter close on the same question. % claim-ok: Theorem 1 of the cited source states the equivalence

The conjecture is open, and this section states it as a problem. % claim-ok: own signpost introducing the moved problem, which carries its own citation

\begin{problem}[{\cite[p.~1]{BelovKanelKontsevich2007}}]\label{prob:open:dixmier}
Let $A_{n,k}$ be the $n$-th Weyl algebra over a field $k$ of characteristic $0$.
Is every algebra endomorphism of $A_{n,k}$ an automorphism?
\end{problem}

Bavula notes that the Dixmier conjecture has been open since 1968 for every
$n\ge1$ \cite[pp.~1--2]{Bavula2017}. Belov-Kanel and Kontsevich prove that the
Jacobian conjecture in dimension $2n$ implies the Dixmier conjecture in
dimension $n$, so the two stable forms are equivalent
\cite[Theorem~1, p.~3]{BelovKanelKontsevich2007}.

The next section builds the same algebra operator by operator, as the
differential operators of the polynomial ring itself
\cite[Theorem~6.2, p.~4]{Hazewinkel2013}.
% claim-ok: own signpost to the next section, moved after the problem

\hbfigure{12-polynomial-weyl/normal-order-rewrite}{The relation $[p,q]=1$ of Definition~\ref{def:pw:weyl} \cite[p.~209]{Dixmier1968}, written out for one pair as $qp=pq-1$. The ordered products $p^{i}q^{j}$ form a basis \cite[p.~210]{Dixmier1968}, so a word is brought to that order by moving each $q$ past the copies of $p$, and every crossing adds the scalar $-1$.}{fig:weyl:normal-order}

Kitchin classifies the derivations of degree-one generalized Weyl algebras over a univariate Laurent polynomial ring, covering the Weyl--Hayashi algebra, a quantisation of the first Weyl algebra arising as a primitive factor algebra of $U_q^+(\mathfrak{so}_5)$, and a family of related algebras~\cite{WlDerivationsOfQuantumAndInvolut2024}.

Brzezi\'nski discusses noncommutative differential geometry of $\mathbb{Z}$-graded generalized Weyl algebras $A(p;q)$ over the polynomial ring in two variables and their zero-degree subalgebras $B(p;q)$, constructs three classes of skew derivations of $A(p;q)$, describes the three-dimensional first-order differential calculi they induce, and computes the associated integrals~\cite{WlNoncommutativeDifferentialGeom2016}.

Bhowmick, Goswami and Landi study covariant derivatives on a class of centred bimodules over an algebra, identify a $Z(A)$-submodule that serves as the analogue of vector fields and prove it is a Lie algebra, show that connections correspond to covariant derivatives on that module, and recover the classical torsion and metric-compatibility formulas~\cite{WlLeviCivitaConnectionsAndVector2020}.

Choi finds all the derivations of the restricted Weyl type nonassociative algebra $WN_{p,s,0}$ on a Laurent extension, continuing the determination of derivations of Weyl type algebras~\cite{WlDerivationsOfARestrictedWeylTy2006}.
